\documentclass[12pt,a4paper]{article}
\usepackage[utf8]{inputenc}
\usepackage[T2A]{fontenc}
\usepackage[russian]{babel}
\usepackage{amsmath}
\usepackage{hyperref}

\title{Руководство пользователя \\ Игра ``Тетрис'' \\ Версия 1.0}

\begin{document}

\maketitle

\section{Введение}

\textbf{«Тетрис»} — одна из самых популярных игр для консоли Brickgame. Нередко и саму консоль называют тетрисом.

\section{Цель игры}

Цель игры — набор очков за построение линий из генерируемых игрой блоков.

\section{Игровой процесс}

\begin{itemize}
\item Очередной блок, сгенерированный игрой, начинает опускаться вниз по игровому полю
\item Блок движется пока не достигнет нижней границы или не столкнется с другим блоком
\item Пользователь может поворачивать фигуры и перемещать их по горизонтали
\item Игрок старается составлять полные ряды из блоков
\end{itemize}

\section{Управление}

\begin{verbatim}
Начало игры       - Enter
Пауза             - P или backspace
Завершение игры   - Esc
Стрелка влево     - перемещение влево
Стрелка вправо    - перемещение вправо
Стрелка вниз      - падение фигуры
Стрелка вверх     - не используется в данной игре
Пробел            - вращение фигуры
\end{verbatim}

\section{Система очков и уровней}

\begin{itemize}
\item За каждый заполненный ряд игрок получает очки
\item После заполнения ряд уничтожается
\item Блоки, находящиеся выше заполненного ряда, опускаются вниз

\item Начисление очков будет происходить следующим образом:
\item 1 линия — 100 очков;
\item 2 линии — 300 очков;
\item 3 линии — 700 очков;
\item 4 линии — 1500 очков.

\item Каждый раз, когда игрок набирает 600 очков, уровень увеличивается на 1
\item Повышение уровня увеличивает скорость движения фигур.
\item Максимальное количество уровней — 10.
\end{itemize}

\section{Завершение игры}

Игра заканчивается, когда очередная фигура останавливается в самом верхнем ряду.

\end{document}